\section{A suíte, e o que ela já pegou}
\label{sec:verificacao}

\status{Pronta. O estado relatado é o da execução de 12 de outubro de 2026: 141
casos, nenhum falhando, contrato 69/69. Os defeitos da seção
\ref{sec:verif-defeitos} são todos reais, e cada um está nomeado com o oráculo que
o encontrou.}

\subsection{Contrato lido pelo arreio}

A suíte é organizada como um \textbf{contrato}: um cabeçalho com cláusulas
numeradas, cada uma declarando uma propriedade que o sistema afirma, explicando
por que ela merece ser afirmada, e nomeando o caso de teste que a verifica. O
arreio de execução \emph{lê as cláusulas do próprio cabeçalho} e exige um caso
passando para cada uma --- cláusula sem teste sai como \textbf{falha}, não como
omissão.

Isso não é formalidade. Em um dos passos deste trabalho, um refator removeu o caso
que cobria uma cláusula, e o arreio reportou ``caso ausente'' nomeando a cláusula
--- o que uma suíte que apenas conta testes passando não faria.

O estado atual, em duas dimensões:

\begin{center}
\begin{tabular}{llr}
\toprule
grupo & assunto & cláusulas \\
\midrule
N & mínimos quadrados móveis: valor, derivada, recusas & 7 \\
B, F & floresta, blocos, suporte & 4 \\
M, R & malha, numeração, refino & 13 \\
G, A, L & GPU, arranjo deslocado, lote de estênceis & 15 \\
E, C & estêncil, condições de contorno por região & 12 \\
P, Q, T & Poisson, compacto, malha adaptada & 11 \\
D & derivada de ordem alta e quadratura de Gauss & 7 \\
\midrule
\multicolumn{2}{l}{\textbf{total}} & \textbf{69} \\
\bottomrule
\end{tabular}
\end{center}

Quatro guardas no arreio merecem menção, porque cada uma existe por um modo de
falha observado: ela nomeia o prefixo de cláusula sem cobertura em vez de apenas
somar; faz \verb|make clean| por configuração, para que um binário compilado com
outro sinalizador não seja medido como se fosse deste; exige a linha de resumo do
executável, para que um processo que morreu antes de terminar não conte como
aprovado; e executa tudo por \verb|mpirun| mesmo com um processo, para que o
caminho de execução seja o mesmo do paralelo.

\subsection{Quatro defeitos que só o contraste encontrou}
\label{sec:verif-defeitos}

Esta subseção existe porque a disciplina de oráculo é barata de descrever e difícil
de levar a sério, e o argumento mais convincente a seu favor é a lista do que ela
pegou neste trabalho.

\paragraph{1. O estêncil medido onde o ajuste é degenerado.} Os primeiros testes
de estêncil e de contorno mediam o ajuste \emph{no centro da célula}. Ali o MLS é
degenerado por construção --- peso $1$ na própria célula, lado direito
$\approx 0$ --- de modo que o fechamento de contorno \textbf{nunca era
exercitado}. Os casos passavam. O que denunciou foi o contraste: ele dava apenas
\num{1e-10} quando devia dar ordens de grandeza. Com alvos deslocados do centro, o
contraste passou a \num{1.9e-1}, e uma cláusula nova agora \emph{afirma} a
degeneração no centro, para que ninguém a redescubra como surpresa.

\paragraph{2. A correção transversal com sinal dependente de quem é grosso.} Na
transição 2:1, fazer o sinal da correção transversal depender de qual lado é o
grosso \textbf{dobrou} o erro: \num{7.8e-2} contra \num{3.9e-2}. A fórmula correta
é independente disso. O sintoma --- erro exatamente duas vezes maior --- foi o que
identificou a causa.

\paragraph{3. O \emph{nan} silencioso do ajuste de ordem 4.} Este é o mais grave,
e é o que motivou o contador de rebaixamento. A rotina de eliminação recusava
pivô somente abaixo de \num{1e-300} --- isto é, apenas a singularidade exata --- e,
ao recusar, devolvia um \textbf{vetor de zeros}. Em malha uniforme isso nunca
ocorria. Em malha adaptada com ordem 4, ocorria em 8 de \num{5172} ajustes, e o
efeito era uma linha da matriz silenciosamente errada:

\begin{center}
\begin{tabular}{lr}
\toprule
nível & $\|e\|_\infty$ (antes do conserto) \\
\midrule
3 & \num{2.36e-10} \\
4 & \textbf{\num{3.32e-1}} \\
5 & \num{2.35e-10} \\
\bottomrule
\end{tabular}
\end{center}

Um nível com erro de \num{3.3e-1} cercado por \num{1e-10} nos vizinhos. O conserto
tem duas partes --- a base adimensionalizada e o rebaixamento contado --- e ambas
estão descritas na seção \ref{sec:aprox-mls}.

\paragraph{4. O ponteiro nulo no ramo que não o usa.} Ao estender a montagem para
o esquema novo, uma edição por substituição de texto inseriu a leitura dos
contadores de MLS \emph{também} no ramo do esquema compacto, onde o objeto
correspondente é um ponteiro nulo. O compilador não acusou, o esquema novo
funcionava, e o compacto --- que passava antes --- quebrava com violação de
segmentação. O experimento que localizou isso em um passo foi trocar a
\emph{ordem} dos casos no arreio: a falha permaneceu no mesmo ponto lógico, o que
descartou de uma vez a hipótese de acumulação de estado e apontou para o caso.

\subsection{O padrão, destilado}

Os quatro têm a mesma forma: \textbf{a afirmação estava certa e o oráculo não
discriminava}. Em três dos quatro, o que denunciou não foi um caso falhando --- foi
um \emph{contraste insuficiente}, isto é, um número que devia ser grande e era
pequeno. Daí a regra de que nenhuma afirmação de exatidão entra na suíte sem o par
que mede onde ela falha.
